\chapter{Motors, transmissions, and speed control}
\label{ch:motors}

A robot must turn a numerical command into physical motion. This requires
more than a motor: a controller generates a command, a driver supplies power,
a transmission adapts the motion, and a sensor measures the result. This
chapter follows that chain for a motor-driven wheel, then closes the loop
by comparing measured speed with desired speed.

\section{From a command to motion}
An electric motor converts electrical energy into mechanical motion. Its
shaft rotates at angular velocity $\omega_m$. A transmission connects that
shaft to a wheel or another mechanism, whose angular velocity $\omega$ may
be quite different. The controller determines how the motor should operate,
but the energy that moves the load comes from a power supply.

\begin{figure}[htbp]
\centering\begin{tikzpicture}[x=1cm,y=1cm,
 block/.style={draw,rounded corners,align=center,minimum height=1cm,font=\small}]
\node[block,text width=2.4cm] (c) at (0,0) {Microcontroller\\controller + PWM};
\node[block,text width=1.4cm] (d) at (3,0) {Motor\\driver};
\node[block,text width=1.1cm] (m) at (5.5,0) {Motor};
\node[block,text width=2.1cm] (g) at (8.3,0) {Transmission};
\node[block] (s) at (3,1.8) {Power supply};
\node[block,text width=1.5cm] (e) at (4,-1.8) {Encoder};
\draw[axis] (c)--node[above=10pt]{command}(d);
\draw[axis] (s)--(d);
\draw[axis] (d)--(m);
\draw[axis] (m)--node[above]{$\omega_m$}(g);
\draw[axis] (g.east)--++(1.1,0) node[right]{$\omega$};
\draw[axis] (6.7,0)|-(e.east);
\draw[axis,sensorframe] (e.west)-|node[pos=.2,below]{counts}(c.south);
\end{tikzpicture}

\caption{The actuation chain. The encoder measures the motor shaft before
the transmission. Its signals return to the microcontroller for speed estimation
and feedback. The motor power supply feeds the driver directly.}
\label{fig:l08-chain}
\end{figure}

The driver, also called a power amplifier, bridges the gap between a small
control signal and the power required by the actuator. A microcontroller's
digital output is intended to communicate a logic state; it is not a motor
power source. Motor current, particularly during starting or stall, can exceed
what a logic pin can supply. The driver switches energy from its supply to
the motor according to the controller's command. The controller and driver
must also have compatible signal levels and, in an ordinary nonisolated
connection, a common electrical reference.

A driver does not create energy or necessarily reproduce the input voltage
amplitude. For example, a logic-level command can control switching of a
separate motor supply at a different voltage. Driver and supply ratings must
match the motor and load; connector size alone does not establish those ratings.

\section{Pulse-width modulation}
\subsection{Representing an intermediate command with two levels}
A digital output can hold a low or a high voltage. To obtain an adjustable
average without generating every intermediate voltage directly, it can switch
periodically between those two levels. In \emph{pulse-width modulation}
(PWM), the duration of the high part of each period sets the command.

Let $T$ be the period and $t_{\mathrm{on}}$ the time spent high. The
\emph{duty cycle} and switching frequency are
\begin{equation}
 D=\frac{t_{\mathrm{on}}}{T},\qquad 0\le D\le1,
 \qquad f_{\mathrm{PWM}}=\frac1T.
\end{equation}
For levels $V_L$ and $V_H$, the average over one period is
\begin{equation}
 \overline V=DV_H+(1-D)V_L.
 \label{eq:l08-average}
\end{equation}
With $V_L=0$ and $V_H=5\,\mathrm V$, a 50\% duty cycle has average
$2.5\,\mathrm V$; a 75\% duty cycle has average $3.75\,\mathrm V$.
The instantaneous voltage still switches between 0 and $5\,\mathrm V$.

\begin{figure}[htbp]
\centering\begin{tikzpicture}[x=1.2cm,y=.6cm]
\draw[axis] (0,0)--(8.6,0) node[right]{$t$ (ms)};
\draw[axis] (0,0)--(0,6) node[above]{$V$};
\foreach \i in {0,...,7} {
 \draw[gray!30] (\i,0)--(\i,5.3);
 \node[below] at (\i,0) {\i};
}
\foreach \i in {0,...,3} {
 \draw[thick] (\i,0)--(\i,5)--(\i+.5,5)--(\i+.5,0)--(\i+1,0);
}
\foreach \i in {4,...,7} {
 \draw[thick] (\i,0)--(\i,5)--(\i+.75,5)--(\i+.75,0)--(\i+1,0);
}
\draw[sensorframe,dashed,thick] (0,2.5)--(4,2.5);
\draw[sensorframe,dashed,thick] (4,3.75)--(8,3.75);
\node[left] at (0,5) {5};
\node[left,sensorframe] at (0,2.5) {2.5};
\node[right,sensorframe] at (8,3.75) {3.75};
\node[above] at (2,5.4) {$D=0.50$};
\node[above] at (6,5.4) {$D=0.75$};
\end{tikzpicture}

\caption{PWM with a fixed period and a change in duty cycle. Dashed blue
lines show the per-period averages, not additional voltage levels produced
by the digital pin. Here $T=1\,\mathrm{ms}$ and $f_{\mathrm{PWM}}=1\,\mathrm{kHz}$.}
\label{fig:l08-pwm}
\end{figure}

\subsection{Why the shaft does not start and stop at every pulse}
The motor and load have inertia, so their speed cannot change instantaneously.
Winding inductance also affects how rapidly current changes. With a suitable
switching frequency, the mechanical response is much smoother than the
electrical waveform. An averaged description is then useful, although
current and torque can still have ripple.

This does not mean that angular velocity equals average voltage: they have
different units. Speed depends on the motor, supply, load, friction, and
driver. Increasing duty cycle generally increases the drive available in a
fixed direction, but a duty cycle alone does not specify an exact speed.
Likewise, zero duty does not necessarily stop a moving wheel immediately;
the result depends on inertia and the driver's braking or coasting behavior.

\subsection{Digital resolution}
An 8-bit command has $2^8=256$ possible values, commonly numbered 0 through
255. In a simple endpoint-inclusive command model,
\[
 D\approx\frac n{255},\qquad n\in\{0,1,\ldots,255\}.
\]
For example, $n=128$ requests approximately 50\% duty. Exact timing depends
on the timer and its configuration. PWM-capable pins, frequency, voltage,
and resolution are hardware-dependent; they are not universal properties of
all microcontrollers.\footnote{For a concrete implementation, see
\href{https://support.arduino.cc/hc/en-us/articles/9350537961500-Use-PWM-output-with-Arduino}{Arduino's PWM documentation}.}

\section{Other forms of actuation}
Electrical motors are common in mobile robots, but the same separation
between command and power applies to fluid-powered actuators. Hydraulic
systems use a liquid; pneumatic systems use a gas. A controlled valve routes
pressurized fluid to an actuator, rather than an electrical driver supplying
motor current.

\begin{figure}[htbp]
\centering\begin{tikzpicture}
\fill[sensorframe!12] (0,0) rectangle (2.7,1.6);
\fill[robotframe!12] (3,0) rectangle (5,1.6);
\draw[thick] (0,0) rectangle (5,1.6);
\filldraw[fill=gray!40] (2.7,0) rectangle (3,1.6);
\filldraw[fill=gray!20] (3,.65) rectangle (6.5,.95);
\node at (1.1,.8) {A: $p_A,A_A$};
\node at (4,1.25) {B: $p_B,A_B$};
\draw[axis,sensorframe] (1.8,.35)--(2.6,.35);
\draw[axis,robotframe] (4,.35)--(3.1,.35);
\draw[axis] (5.5,1.3)--(6.5,1.3) node[right]{positive motion};
\draw (1,0)--(1,-.5) node[below]{fluid port};
\draw (4,0)--(4,-.5) node[below]{fluid port};
\end{tikzpicture}

\caption{A double-acting cylinder, shown schematically. Fluid pressure acts
on the two faces of the piston. The rod makes their effective areas unequal.}
\end{figure}

For the cylinder shown, pressure in chamber A pushes the piston to the right,
while pressure in chamber B pushes it to the left. Using pressures relative
to the surrounding pressure, the ideal fluid force to the right is
\[
 F=p_A A_A-p_B A_B.
\]
External load and friction also enter the force balance. Motion therefore
does not simply continue until the chamber pressures are equal: the effective
areas and applied load matter. The remaining sections focus on a rotary
electric motor driving a wheel.

\section{Transmissions and reduction ratio}
A motor may rotate faster than the wheel needs to rotate. A gear train can
reduce this speed and adapt the available torque. Transmissions can also
relocate an axis, change its direction, or convert rotation into translation.

Define the positive reduction ratio by the speed magnitudes,
\begin{equation}
 R=\frac{|\omega_m|}{|\omega|}.
 \label{eq:l08-ratio}
\end{equation}
For a reduction, $R>1$. Once the positive directions of the two shafts have
been chosen to correspond through the transmission, the signed relation is
\begin{equation}
 \omega=\frac{\omega_m}{R}.
 \label{eq:l08-wheelomega}
\end{equation}
This convention absorbs any reversal caused by the gear train into the shaft
coordinate definitions. If both axes instead use the same physical viewing
direction, a reversing transmission requires a minus sign.

\begin{figure}[htbp]
\centering\begin{tikzpicture}[box/.style={draw,minimum height=1.1cm,align=center}]
\node[box] (e) at (0,0) {Encoder\\$N$ ticks/rev};
\node[box,minimum width=1.6cm] (m) at (2.4,0) {Motor};
\node[box,minimum width=2cm] (g) at (5.3,0) {Gearbox\\ratio $R$};
\draw[thick] (e)--(m);
\draw[thick] (m)--node[above]{$\omega_m$}(g);
\draw[thick] (g.east)--(8.2,0);
\draw[thick] (8.2,0) circle (1);
\fill (8.2,0) circle (.04);
\draw[axis,sensorframe] (8.2,0)--node[right]{$r$}(8.2,-1);
\draw[axis,robotframe] (8.2,1.3) arc (90:-60:1.3);
\node[robotframe] at (9.7,.8) {$\omega$};
\draw[thick] (6.8,-1)--(9.7,-1);
\draw[axis] (7.4,-1.7)--(9,-1.7) node[right]{$v=r\omega$};
\node at (8.2,1.8) {Wheel};
\end{tikzpicture}

\caption{Motor-side sensing and wheel-side motion. The gearbox represents
the entire transmission, not a particular number or arrangement of gears.
The encoder resolution refers to the shaft on which it is mounted.}
\label{fig:l08-transmission}
\end{figure}

For $R=100$, one wheel revolution requires 100 motor revolutions. Reduction
does not provide free mechanical power: ideally, the decrease in speed is
accompanied by an increase in torque, while real transmissions dissipate
some energy. Backlash and compliance can also make wheel motion differ
slightly from the motion inferred at the motor shaft.

\section{Encoders and speed measurement}
An incremental encoder produces transitions as its shaft turns. A magnetic
encoder can use a rotating magnet and Hall-effect sensors; an optical encoder
uses a different sensing principle to obtain similar position increments.
Counting these increments over a known time interval gives a speed estimate.

One periodic channel provides information about the amount of rotation but
does not, by itself, distinguish the two directions. Two suitably phased
channels allow direction to be inferred from their transition order.
In a quadrature encoder, channels A and B are offset by a quarter of an
electrical cycle. Reversing rotation reverses which channel leads. This phase
offset need not equal a 90-degree mechanical separation of the sensors.

\begin{figure}[htbp]
\centering\begin{tikzpicture}[x=1cm,y=.7cm]
\node[left] at (0,2.5) {A};
\node[left] at (0,.5) {B};
\draw[gray!40] (0,2)--(8.5,2);
\draw[gray!40] (0,0)--(8.5,0);
\draw[thick,sensorframe] (0,2)--(1,2)--(1,3)--(3,3)--(3,2)--(5,2)--(5,3)--(7,3)--(7,2)--(8.5,2);
\draw[thick,robotframe] (0,0)--(2,0)--(2,1)--(4,1)--(4,0)--(6,0)--(6,1)--(8,1)--(8,0)--(8.5,0);
\foreach \i in {1,...,8} {\draw[projection,gray] (\i,-.3)--(\i,3.3);}
\draw[axis] (0,-.7)--(9,-.7) node[right]{$t$};
\draw[<->] (1,3.7)--node[above]{one electrical cycle}(5,3.7);
\end{tikzpicture}

\caption{Ideal quadrature signals for one direction. Channel A rises before
channel B. Reversing the direction reverses the sequence. Counting all edges
of both channels gives four counts per electrical cycle.}
\end{figure}

\subsection{Keeping the counting convention explicit}
Let $N$ denote the number of \emph{counted events per motor revolution},
also called ticks per revolution (TPR) here. It must match the decoding
method: counting rising edges of one channel is different from counting
both edges of both channels. Manufacturer terms such as pulses and counts
must be interpreted together with the stated convention.

If $\Delta n$ signed counts are accumulated over an interval $\Delta t$, then
\begin{equation}
 \overline\omega_m=\frac{2\pi\Delta n}{N\Delta t}
 \quad\text{in radians per second}.
 \label{eq:l08-encoder}
\end{equation}
The overbar denotes an estimate obtained from the encoder. In particular,
the count rate $q=\Delta n/\Delta t$ in ticks per second is converted to
angular velocity by multiplying by $2\pi/N$.

An estimate is not the exact instantaneous speed. Finite counts cause
quantization, and the counting window averages motion over time. A longer
window can improve count resolution while delaying the response to speed
changes. Missed transitions and mechanical uncertainty introduce further
errors. At low speed, measuring the time between transitions is another
possible approach.

\section{From encoder counts to wheel speed}
For a wheel of radius $r$ rolling without slip along a stationary surface,
its center's velocity component in the rolling direction satisfies
\begin{equation}
 v=r\omega.
 \label{eq:l08-rolling}
\end{equation}
The same quantity is the magnitude of the rim's tangential velocity relative
to the wheel center. It is not the ground-frame velocity of the material
point touching the ground: that point is instantaneously at rest in ideal
rolling. For a robot with multiple wheels, their velocities must be combined
using the robot's kinematic model to obtain its overall motion.

Combining the encoder conversion, reduction, and rolling relation gives
\begin{equation}
 \boxed{\overline v=\frac{2\pi r}{NR}\frac{\Delta n}{\Delta t}.}
 \label{eq:l08-speed}
\end{equation}
This is a wheel-based estimate. If the wheel slips, encoder measurements can
still describe shaft rotation accurately while failing to describe travel
over the ground.

\subsection{Worked example: reconstructing motion}
Suppose the motor encoder provides $N=8$ counted ticks per revolution,
the reduction is $R=100$, and the wheel radius is $r=0.05\,\mathrm m$.
The measured count rate is $q=40\,\mathrm{ticks/s}$. First convert units:
\begin{align*}
 \overline\omega_m
 &=40\,\frac{\mathrm{ticks}}{\mathrm s}
   \frac{1\,\mathrm{rev}}{8\,\mathrm{ticks}}
   \frac{2\pi\,\mathrm{rad}}{\mathrm{rev}}\\
 &=10\pi\,\mathrm{rad/s}\approx31.416\,\mathrm{rad/s}.
\end{align*}
The transmission then gives
\[
 \overline\omega=\frac{10\pi}{100}
 =0.1\pi\,\mathrm{rad/s}\approx0.31416\,\mathrm{rad/s}.
\]
Finally,
\[
 \overline v=0.05(0.1\pi)
 =0.005\pi\,\mathrm{m/s}\approx0.015708\,\mathrm{m/s}.
\]
Thus the estimated rolling speed is about $1.57\,\mathrm{cm/s}$.
The first step changes units; the second describes the transmission; the
third converts angular motion into linear motion using the wheel radius.

\section{From desired wheel speed to desired motor speed}
The same relationships can be used in reverse. A desired rolling speed
$v_d$ implies
\begin{equation}
 \omega_d=\frac{v_d}{r},\qquad
 \boxed{\omega_{m,d}=\frac{R}{r}v_d.}
 \label{eq:l08-desired}
\end{equation}
Using the same hardware and $v_d=0.1\,\mathrm{m/s}$,
\[
 \omega_d=\frac{0.1}{0.05}=2\,\mathrm{rad/s},\qquad
 \omega_{m,d}=100(2)=200\,\mathrm{rad/s}.
\]
The corresponding expected encoder count rate would be
\[
 q_d=\frac{N}{2\pi}\omega_{m,d}
 =\frac{800}{\pi}\,\mathrm{ticks/s}
 \approx254.65\,\mathrm{ticks/s}.
\]
This calculation specifies a target motor speed. It does not yet determine
which PWM command will achieve it under a particular load.

\section{Closing the speed-control loop}
In open-loop operation, a chosen command is applied without correcting it
from the measured outcome. Feedback instead compares the desired motor
speed with the measured motor speed:
\begin{equation}
 e(t)=\omega_{m,d}(t)-\overline\omega_m(t).
 \label{eq:l08-error}
\end{equation}
Both terms must use the same units and shaft reference. Subtracting a wheel
speed from a motor speed, or radians per second from ticks per second,
would not produce the intended error.

For a fixed forward direction, a positive error indicates that the motor is
too slow; a negative error indicates that it is too fast. A controller uses
this information to adjust the command. The encoder return path in
Figure~\ref{fig:l08-chain} completes this feedback loop. A motor operated
with such feedback is part of a servo system; this does not imply that it is
an integrated hobby position servo.

\subsection{Setting a command versus changing a command}
Let $u$ be a normalized drive command. A proportional controller directly
sets the command from the current error:
\begin{equation}
 u(t)=K_P e(t),\qquad K_P>0.
 \label{eq:l08-p}
\end{equation}
For a forward-only drive, the applied duty is limited to $[0,1]$. A driver
that supports reversal can instead interpret a signed command through
separate direction and duty settings.

There is a useful distinction between Equation~\eqref{eq:l08-p} and the
instruction ``increase the existing command while the error is positive.''
If the latter is implemented as
\begin{equation}
 \dot u(t)=k e(t),
 \qquad\text{then}\qquad
 u(t)=u(0)+k\int_0^t e(\tau)\,d\tau.
 \label{eq:l08-incremental}
\end{equation}
It accumulates error and is integral action on the command. With sample
interval $T_s$, a corresponding update is
$u_{k+1}=u_k+kT_s e_k$, followed by the appropriate command limits.
The distinction matters: a dot over the command changes the controller.

\subsection{A brief introduction to PID}
A conventional ideal proportional--integral--derivative controller combines
three contributions to the command itself:
\begin{equation}
 u(t)=K_P e(t)+K_I\int_0^t e(\tau)\,d\tau
       +K_D\frac{de(t)}{dt}.
 \label{eq:l08-pid}
\end{equation}
The proportional term responds to present error, the integral term to
accumulated error, and the derivative term to the rate of change of error.
The gains have different units because these three inputs have different
units.\footnote{See the parallel PID form in
\href{https://www.mathworks.com/help/control/ug/proportional-integral-derivative-pid-controllers.html}{MathWorks' PID overview}.}

Proportional control can leave a steady speed error when a nonzero command
is needed to sustain motion. Integral action can remove a constant offset
when the closed loop is stable and the requested speed is achievable.
Derivative action is sensitive to measurement noise and is usually filtered
in practical implementations. Command saturation also needs attention:
continued accumulation while the drive is already at its limit can produce
integral windup.

Choosing gains determines how quickly the system responds and whether it
oscillates or becomes unstable. A fast processor alone does not guarantee
good control. The feedback update period and the PWM switching period are
distinct: many switching cycles may occur between feedback updates. Detailed
controller design is beyond the present introduction; the essential idea is
to measure the outcome and use the error to correct the drive.

\clearpage
\section{Additional exercises}
Use positive directions that correspond through the transmission. Unless
otherwise stated, the encoder is on the motor shaft and the wheel rolls
without slip.
\begin{enumerate}
\item \textbf{PWM timing.} A signal switches between 0 and $5\,\mathrm V$
at $1\,\mathrm{kHz}$. Find its period, high time, and average voltage for
duty cycles of 20\%, 50\%, and 80\%. Explain why the average voltage is
not an angular velocity.

\item \textbf{Command resolution.} Use $D=n/255$ for an 8-bit command.
Choose the nearest integer command for 60\% duty and calculate the resulting
average of a 0--$5\,\mathrm V$ signal. How many command values are available?

\item \textbf{Encoder conversion.} A motor encoder produces 12 counted
events per revolution. During $0.2\,\mathrm s$, the controller records 36
forward counts. Find the count rate, motor speed in revolutions per second,
and motor speed in radians per second. For $R=50$ and $r=0.04\,\mathrm m$,
find the wheel's rolling speed.

\item \textbf{Desired motion.} For the hardware in Exercise 3, determine
the desired motor speed and expected count rate for $v_d=0.12\,\mathrm{m/s}$.
Why does this not determine the required duty cycle uniquely?

\item \textbf{Counting conventions.} A quadrature encoder produces eight
complete electrical cycles per motor revolution on each channel. Determine
$N$ when counting only rising edges of A and when counting every edge of
both A and B. What speed error results if software uses the first value of
$N$ while actually counting with the second method?

\item \textbf{Resolution and slip.} Using $N=8$, $R=100$, and
$r=0.05\,\mathrm m$, find the inferred travel per count. Find the change
in estimated speed caused by one extra count in windows of $0.01$ and
$0.1\,\mathrm s$. Explain the resolution--delay tradeoff and why neither
window can correct wheel slip by itself.

\item \textbf{Feedback.} Let $\omega_{m,d}=200\,\mathrm{rad/s}$ and
$\overline\omega_m=170\,\mathrm{rad/s}$. Compute the error. With
$K_P=0.01\,\mathrm{s/rad}$, find the proportional command $u=K_Pe$.
Separately, starting from $u_k=0.4$, evaluate the incremental update
$u_{k+1}=u_k+kT_s e_k$ for $k=0.2\,\mathrm{rad}^{-1}$ and
$T_s=0.01\,\mathrm s$. Explain why these are different controllers.
\end{enumerate}
